\section{Computation of local invariants}
\label{sec-computation-local-invariants}

\subsection{General considerations and statement}

Our goal in this section is to compute the correlated GW invariants in the case where $X=E$ is an elliptic curve, the genus of the source curve is $g=1$, and with a unique interior point constraint. These \textit{local invariants} will be used in combination with Theorem~\ref{theo-refined-decomposition-surjective-case} to obtain an explicit computation algorithm and derive regularity results for general correlated invariants in Section \ref{sec-regularity}.

We consider $Y=\PP(\O\oplus L)$ for $L\in\Pic^0(E)$; fix an homology class $\beta=a[E]\in H_2(E,\ZZ)$ and a vector $\bfw=(w_1,\dots,w_n)$ of tangency orders. The moduli space
\[ \M(a,\bfw)=\M_{1,1}\bigl(Y|D^\pm,a[E],\bfw\bigr) \]
is the moduli space of log stable maps as in the previous section.
For $\delta|\mathrm{gcd}(w_i)$, we saw that it decomposes into components $\M^\theta(a,\bfw)$.
The correlators $\theta$ satisfy
$
\delta\theta\equiv \varphi_{a[E]}(L)=a\lambda$,
where~\smash{$\lambda=\varphi_{[E]}(L)$} is the image of the line bundle $L\in\Pic^0(E)\simeq E$ through the isomorphism~$\varphi_{[E]}$. We denote the torsor by $T^\delta(L,a)$.

\subsubsection{Uncorrelated case} We recall the computation of the non-refined invariant
\[ \langle \pt_0,1_{w_1},\pt_{w_2},\dots,\pt_{w_n}\rangle_{1,a[E],\bfw}=
\int_{\vir{\M(a,\bfw)}}\ev_0^*(\pt)\prod_2^n\ev_i^*(\pt). \]

\begin{lem}\label{lem-non-refined-computation}
The uncorrelated relative GW invariant has the following value:
\[ \langle \pt_0,1_{w_1},\pt_{w_2},\dots,\pt_{w_n}\rangle_{1,a[E],\bfw}=a^{n-1}\sigma(a)\cdot w_1^2, \]
where $\sigma(a)=\sum_{d|a}d$ is the sum of divisors function.
\end{lem}

\begin{proof}
This elementary formula is computed, for instance, in \cite{blomme2022floor}. Briefly, it may be obtained as follows. Genus $1$ parametrized curves in $Y$ realizing the class $a[E]+b(\bfw)\big[\PP^1\big]$ come from degree~$a$ covering maps $f\colon C\to E$. Up to translation, these are group homomorphisms. We~may fix the translation parameter using the marked point $p_0$, mapped to a fixed point in $E$.
Then we see that these are in bijection with the index $a$ sublattices of $\pi_1(E)\simeq\ZZ^2$, of which there are~$\sigma(a)$.

 Given one of the $\sigma(a)$ covers $f\colon C\to E$, enhancing to a map to $Y=\PP(\O\oplus L)$ amounts to find a section of $f^*L$, with poles and zeros of prescribed order $w_i$ and fixed image in $E$.

We denote by $y_i$ the %position of
marked points in $C$ (with not trivial contact order) and let $x_i=f(y_i)$ their image in $E$. For $2\leqslant i\leqslant n$, there are $a$ possible choices for each $y_i$ given that $x_i$ is fixed. The position of $y_1$ is determined by the relation $\sum w_iy_i\equiv f^*(\lambda)$ in $C$, i.e., $y_1$ is a $w_1$-root of~${f^*(\lambda)-\sum_{i=2}^nw_iy_i}$.
There are $w_1^2$ possible choices for $y_1$, yielding the result.
\end{proof}

\subsubsection{Statement} We now consider the correlated invariants
\[ \langle\langle \pt_0,1_{w_1},\pt_{w_2},\dots,\pt_{w_n}\rangle\rangle^\delta_{1,\beta,\bfw}
= \int_{[\![\M(a,\bfw)]\!]^\delta}\ev_0^*(\pt)\prod_2^n\ev_i^*(\pt), \]
which is an element of the group algebra $\QQ[E]$ with support on $T^\delta(L,a)$.



Before giving a closed formula for the full local correlated invariant, we need to introduce certain functions with values in $\QQ[E]$.

First, we consider the average of torsion elements
\[\vartheta_d:=\frac{1}{d^2}\sum_{d\theta\equiv 0}(\theta) \in\QQ[E].\]
They satisfy \smash{$\vartheta_{d_1}\vartheta_{d_2}=\vartheta_{\mathrm{lcm}(d_1,d_2)}$}.
We then define the following functions: for $d|\delta$,
\[
 \boldsymbol{\vartheta}_\delta(d) = \prod_p (\vartheta_{p^{\nu_p(d)}}-\mathds{1}_{\nu_p(d)<\nu_p(\delta)}\vartheta_{p^{\nu_p(d)+1}} ),
 \]
where the products is over primes and $\nu_p$ is the $p$-adic valuation. The $\mathds{1}$ indicates that the second term for each factor of the product only appears if $\nu_p(d)<\nu_p(\delta)$.

\begin{expl}
 If $d=\delta$, so that $\boldsymbol{\vartheta}_\delta(\delta) =\vartheta_\delta$. If $\delta=p^v$ is the power of a prime number $p$, the values of $\boldsymbol{\vartheta}_{p^v}$ for its divisors $1,p,\dots,p^{v-1},p^v$ are
 $\vartheta_1-\vartheta_p,\vartheta_p-\vartheta_{p^2},\dots, \vartheta_{p^{v-1}}-\vartheta_{p^v}$ and $\vartheta_{p^v}$.
\end{expl}

Next, we consider the variants of the arithmetic function $\sigma(a)=\sum_{d| a}d$ defined as follows: \smash{$\overline{\sigma}^d(a)=\sigma(a/d)$} if $d|a$ and $0$ else. Combining both, we define the following function, with values in $\QQ[\Tor_\delta(E)]$
\[ \boldsymbol{\sigma}_\delta(a)=\sum_{d|\delta}\overline{\sigma}^{\delta/d}(a)\boldsymbol{\vartheta}_\delta(d). \]
For concrete computations, it may be convenient to express $\boldsymbol{\sigma}_\delta(a)$ as a linear combination of the~$\vartheta_d$ for $d|\delta$ but with different coefficients. To do so, consider the following function on $\NN$:
\[
\Upsilon^\delta_d(a) = \prod_p \bigl(\overline{\sigma}^{p^{\nu_p(d)}}-\mathds{1}_{\nu_p(d)<\nu_p(\delta)}\overline{\sigma}^{p^{\nu_p(d)+1}} \bigr)\bigl(p^{\nu_p(a)}\bigr).
\]
In the above product of function, the argument $a$ is factorized over prime numbers, so that for~$a=\prod p^{\nu_p(a)}$ and $d=\prod p^{\nu_p(d)}$,
\[ \overline{\sigma}^d(a)=\prod\overline{\sigma}^{p^{\nu_p(d)}}\bigl(p^{\nu_p(a)}\bigr). \]
The following identity is proven in the proof of Theorem~\ref{theo-expression-local-correlated-invariant}. Thanks to multiplicativity, it suffices to check it for powers of primes, where it stems from a summation by part
\[ \boldsymbol{\sigma}_\delta(a)=\sum_{d|\delta}\Upsilon_d^\delta(a)\vartheta_{\delta/d}. \]
Doing a summation by parts for each prime numbers is called a \textit{multiplicative summation by parts}.

The torsor consists of $\delta$-roots of $a\lambda$. However, $a\lambda$ has no canonical $\delta$-root. The furthest we can naturally define is the $\gcd(a,\delta)$-root \smash{$\frac{a\lambda}{\gcd(a,\delta)}$}. Then, pick $\theta_0$ to be any choice of root such that $\frac{\delta}{\gcd(a,\delta)}\theta_0 = \frac{a}{\gcd(a,\delta)}\lambda$. A correlator $\theta_0$ as before is called a \textit{special correlator}. Theorem~\ref{theo-torsor-invariance} ensures that the result does not depend on the choice of the latter.

One way to construct special correlators is as follows: if $\lambda_0$ satisfies $\delta\lambda_0=\lambda$, one may take~${\theta_0=a\lambda_0}$. Indeed, one has
\[ \frac{\delta}{\gcd(a,\delta)}(a\lambda_0) = \frac{a}{\gcd(a,\delta)}\delta\lambda_0 = \frac{a}{\gcd(a,\delta)}\lambda. \]

\begin{theo}\label{theo-expression-local-correlated-invariant}
The full local correlated invariant has the following expression:
\begin{equation}\label{eq-local-invariant}
\langle\langle \pt_0,1_{w_1},\pt_{w_2},\dots,\pt_{w_n}\rangle\rangle^\delta_{1,a[E],\bfw} = a^{n-1}w_1^2 \boldsymbol{\sigma}_\delta(a)
\cdot (\theta_0),
\end{equation}
where $\theta_0$ is any special correlator, i.e., satisfies \smash{$\frac{\delta}{\gcd(a,\delta)}\theta_0 = \frac{a}{\gcd(a,\delta)}\lambda$}.
\end{theo}

Without assuming Theorem~\ref{theo-expression-local-correlated-invariant}, it can be checked from the definition of $\boldsymbol{\sigma}$ that
\begin{equation}\label{eq-sigma-invariance}
 \boldsymbol{\sigma}_\delta(a) = \boldsymbol{\sigma}_\delta(a)\vartheta_{\delta/\gcd(a,\delta)},
\end{equation}
so that the right-hand side of \eqref{eq-local-invariant} does not depend on which $\theta_0$ we choose. It may thus be replaced by \smash{$\vartheta_{\delta/\gcd(a,\delta)}\cdot(\theta_0)=\dfk\big[\frac{1}{\delta/\gcd(a,\delta)}\big]\bigl(\frac{a}{\gcd(a,\delta)}\lambda\bigr)$}.

Assuming instead Theorem~\ref{theo-expression-local-correlated-invariant}, equation \eqref{eq-sigma-invariance} is in fact a consequence from Theorem~\ref{theo-torsor-invariance}, which tells us that the correlated invariant is invariant when multiplying by an element of \[\varphi_{a[E]}(\Tor_\delta(E))=a\cdot\Tor_\delta(E)=\Tor_{a/\gcd(a,\delta)}(E).\]
See the second part of the proof of Lemma \ref{lem-k=f(tor)} for a proof of the second equality. In particular, it is also invariant by multiplication by $\vartheta_{\delta/\gcd(a,\delta)}$.

With the above formulation, the correlated invariant appears as a refined version of the uncorrelated invariant, and the refinement takes the form of $\boldsymbol{\sigma}_\delta$ replacing $\sigma$.


\subsubsection{Applications} Before going to the proof, we present some applications.

\begin{expl}
 If $a$ is coprime with $\delta$, every $\overline{\sigma}^{\delta/d}(a)$ vanishes except for $d=\delta$, and we thus recover that
$\boldsymbol{\sigma}_\delta(a) = \sigma(a)\vartheta_\delta$.
 In other words, the curves are equally spread among the correlators.
\end{expl}

\begin{expl}
 Let $J_2$ be the second Jordan function, defined by $J_2(p^\alpha)=p^{2\alpha-2}\bigl(p^2-1\bigr)$. It counts the number of elements of order $n$ in $(\ZZ/n\ZZ)^2$.

 The $(0)$-coefficient of $\boldsymbol{\vartheta}_\delta(d)$ is equal to the product of $(0)$-coefficients for each prime $p$. The~latter is equal to $\frac{1}{p^{2\nu_p(d)}}\!-\!\frac{1}{p^{2\nu_p(d)+2}}$ if $\nu_p(d)\!<\!\nu_p(\delta)$ and $\frac{1}{p^{2\nu_p(\delta)}}$ else. This matches the values of~the~function \smash{$\frac{J_2(\delta/d)}{\delta^2}$}, which is also multiplicative. Therefore,
 we get that
 \[ \langle \pt_0,1_{w_1},\pt_{w_2},\dots,\pt_{w_n}\rangle^{\theta_0}_{1,a[E],\bfw} =a^{n-1}\left(\frac{w_1}{\delta}\right)^2\sum_{d|\delta}J_2(d)\overline{\sigma}^d(a). \]
 However, this application is partially a lie since this computation is actually the first step toward proving Theorem~\ref{theo-expression-local-correlated-invariant}.
\end{expl}

From $\m{\delta/\delta'}\bigl(\fvir{\M(a,\bfw)}{\delta}\bigr) = \fvir{\M(a,\bfw)}{\delta'}$, we deduce that the functions $\boldsymbol{\sigma}_\delta$ satisfy
\[ \m{\delta/\delta'}(\boldsymbol{\sigma}_\delta(a)) = \boldsymbol{\sigma}_{\delta'}(a),\]
which may also be checked directly from the definition of $\boldsymbol{\sigma}_\delta$.

Using Theorem~\ref{theo-expression-local-correlated-invariant}, we immediately get the quasi-modularity result for the generating series of correlated invariants. Here, we extend the notion of quasi-modularity for functions with value in a vector space, here chosen to be $\CC[E]$. In the finite-dimensional case, it just means that all coordinate functions are quasi-modular forms. See Section \ref{sec-modularity} for more details.

\begin{coro}
The following generating series is quasi-modular for $\Gamma_0(\delta)$
\[ \sum_a \langle\langle \pt_0,1_{w_1},\pt_{w_2},\dots,\pt_{w_n} \rangle\rangle^\delta_{1,a[E],\bfw}\sfq^a.\]
\end{coro}

\begin{proof}
 The generating series of $\sigma$ is the first Eisenstein series $E_2(\sfq)$, known to be a quasi-modular form. The result thus follows from the quasi-modularity of the generating series $\sum_a \overline{\sigma}^d(a)\sfq^a = E_2(\sfq^d)$ for the congruence subgroup $\Gamma_0(d)$.
\end{proof}

The rest of the section is dedicated to the proof of Theorem~\ref{theo-expression-local-correlated-invariant}, by refining the proof of Lemma \ref{lem-non-refined-computation}. We proceed in several steps. The first is to study the curves coming from a~common covering map $f\colon C\to E$. Summing over covering maps in Section \ref{sec-0-coeff}, we are then able to find a closed expression for the $(\theta_0)$-coefficient. Multiplicativity properties and an induction relation are thus sufficient to prove the formula from Theorem~\ref{theo-expression-local-correlated-invariant}.

 \subsection{Contribution of a fixed cover}

 Let us consider a fixed cover $f\colon C\to E$, which is a group homomorphism choosing the marked point $p_0$ and its image as neutral element. Its kernel $\ker f$ has cardinality $a$, the degree of the covering. Let $f^*\colon E\to C$ be the dual map between the curves seen as their Picard groups. The lifting condition writes itself $\sum_1^n w_jy_j=f^*(\lambda)$, where $\lambda\in E$ corresponds to the line bundle $L$. We have the following group morphism:
 \[
 C^n \longrightarrow E^{n-1}\times C, \qquad
 (k_j) \longmapsto \left(f(k_2),\dots,f(k_n),\sum_1^n w_jk_j \right)
 \]
 with its kernel $K(f)=\big\{ (k_j)\in C\times(\ker f)^{n-1} \text{ s.t. } \sum_1^n w_jk_j=0 \big\}$. In particular, it sits in the following exact sequence, from which we see it has cardinality $w_1^2 a^{n-1}$
 \[ 0\to \Tor_{w_1}(C) \to K(f) \to (\ker f)^{n-1} \to 0. \]

 For the cover $f$, the set of curves matching the constraints $(x_j)$ is in bijection with the following $K(f)$-torsor
 \[ S_{\lambda,x_2,\dots,x_n}=\left\{ (y_j)\in C^n \text{ s.t. }\forall \ 2\leqslant j\leqslant n \ f(y_j)=x_j \text{ and }\sum_1^n w_jy_j=f^*(\lambda)\in C\right\}, \]
also having cardinality $w_1^2a^{n-1}$. We now wish to refine the above description. To do so, we use the correlator function
 \[
 \kappa^\delta \colon\ C^n \longrightarrow E,\qquad
 (y_j) \longmapsto \sum_1^n \frac{w_j}{\delta}f(y_j).
 \]
 The correlators are the elements $\theta\in E$ satisfying $\delta\theta=f_*f^*(\lambda)=a\lambda$. Among them, recall we have the \textit{special correlators} satisfying $\frac{\delta}{\gcd(a,\delta)}\theta_0=\frac{a}{\gcd(a,\delta)}\lambda$ and that $a\lambda_0$, where $\delta\lambda_0=\lambda$ is a~special correlator.

\begin{lem}\label{lem-k=f(tor)}
 The image of $K(f)$ via the correlator function is $\kappa^\delta(K(f)) = f(\Tor_\delta(C))$. Furthermore, it contains \smash{$\Tor_{\delta/\gcd(a,\delta)}(E)$}.
\end{lem}

\begin{proof}
 The equality follows from the definitions. Assume that $(k_j)\in K(f)$. Then we have that $\sum_1^n w_jk_j=0$ and it follows that $\sum_1^n\frac{w_j}{\delta}k_j\in\Tor_\delta(C)$ and consequently $\kappa^\delta((k_j))=f\bigl(\sum_1^n\frac{w_j}{\delta}k_j\bigr)\in f(\Tor_\delta(C))$. Conversely, let $f(s)\in f(\Tor_\delta(C))$ with $s\in\Tor_\delta(C)$. Let $k_1$ be such that $s=\frac{w_1}{\delta}k_1$, which exists because $C$ is divisible, and $k_j=0\in\ker f$ if $j\geqslant 2$. We have
 \[ f(s)=f\left(\frac{w_1}{\delta}k_1\right)=\kappa^\delta(k_1,0,\dots,0)\qquad \text{and}\qquad\sum_1^n w_jk_j=\delta s=0. \]
 Thus, we have the reverse inclusion $f(\Tor_\delta(C)\subset\kappa^\delta(K(f))$.

 We now prove that $\Tor_{\delta/\gcd(a,\delta)}(E)\subset f(\Tor_\delta(C))$. Using the dual morphism $f^*$, we start from $f^*(\Tor_\delta(E))\subset\Tor_\delta(C)$.
 We now apply the morphism $f$ and use that the composition $f\circ f^*\colon E\to E$ is the multiplication by $a$, so that $a\cdot\Tor_\delta(E)\subset f(\Tor_\delta(C))$.
 We now claim that~${a\cdot\Tor_\delta(E)=\Tor_{\delta/\gcd(a,\delta)}(E)}$, which concludes the proof:
 \begin{itemize}\itemsep=0pt %[label=$\triangleright$]
 \item If $x$ is $\delta$-torsion, we have \smash{$\frac{\delta}{\gcd(a,\delta)}ax=\frac{a}{\gcd(a,\delta)}\delta x=0$}, so that we have the inclusion
 \[
 a\cdot \smash{\Tor_\delta(E)\subset\Tor_{\delta/\gcd(a,\delta)}(E)}.
 \]
 \item As $|\Tor_u(E)|=u^2$, the short exact sequence
 \[ 0\to \Tor_a(E)\cap\Tor_\delta(E)=\Tor_{\gcd(a,\delta)}(E) \to \Tor_\delta(E) \xrightarrow{a\cdot} a\cdot\Tor_\delta(E) \to 0, \]
 ensures that they have the same cardinality $\bigl(\frac{\delta}{\gcd(a,\delta)}\bigr)^2$.\hfill $\qed$
 \end{itemize}\renewcommand{\qed}{}
\end{proof}

The following proposition gives a description of the correlators in the image of the torsor~$S_{\lambda,x_2,\dots,x_n}$.

\begin{prop}\label{prop-solution-torsor}
The correlators achieved by the solutions, i.e., the set $\kappa^{\delta}(S_{\lambda,x_2,\dots,x_n})$, form a $f(\Tor_\delta(C))$-torsor which contains the special correlators $\theta_0$. Solutions are uniformely spread among the correlators.
\end{prop}

\begin{proof}
 The solutions form the $K(f)$-torsor $S_{\lambda,x_2,\dots,x_n}$. Applying the correlator function, we immediately get a torsor under $\kappa^\delta(K(f))$, equal to $f(\Tor_\delta(C))$ by Lemma \ref{lem-k=f(tor)}. We also get that the solutions split evenly among the elements.

 We now need to prove that it contains the special correlators. Special correlators form a~torsor under $\Tor_{\delta/\gcd(a,\delta)}(E)$, which lies in $f(\Tor_\delta(C))$ by Lemma \ref{lem-k=f(tor)}. Thus, it suffices to show it contains one of them.
 Pick $y_2,\dots,y_n$ such that $f(y_j)=x_j$ and choose $y_1$ such that~${\sum_1^n\frac{w_j}{\delta}y_j=f^*(\lambda_0)}$, with $\delta\lambda_0=\lambda$; notice that this always exists as $C$ as well is a divisible group. In particular, as $\delta\lambda_0=\lambda$, we have that $\sum_1^n w_jy_j=f^*(\lambda)$ and $(y_j)$ is indeed an element of $S_{\lambda,x_2,\dots,x_n}$. Moreover, we have that
$\kappa^\delta((y_j))=f\circ f^*(\lambda_0)=a\lambda_0$,
 which is one of the special correlators. Therefore, all special correlators belong to the image torsor.
\end{proof}

To finish this section, we provide a first expression of the local correlated invariant as a sum over the covers $f\colon C\to E$. The cover $f$ corresponds to a sublattice $\Lambda\subset\pi_1(E)\simeq\ZZ^2$. We can find a basis $(e_1,e_2)$ such that $\Lambda=\langle ke_1,(a/k)e_2\rangle$ for some unique $k$ such that $k^2|a$. We say that~$\Lambda$, and by extension the associated cover, is of type $(k,a/k)$.

Let $\vartheta(f)\in\ZZ[E]$ be the element with coefficient $1$ for every element in $f(\Tor_\delta(C))$.

\begin{prop}\label{prop-first-expression-local-correlated-inv}
 The correlated invariant admits the following expression:
 \[ \langle\langle\pt_0,1_{w_1},\pt_{w_2},\dots,\pt_{w_n} \rangle\rangle^\delta_{1,a[E],\bfw} = a^{n-1}\left(\frac{w_1}{\delta}\right)^2\sum_{f\colon C\to E} \gcd(k,\delta)\gcd(a/k,\delta) \vartheta(f)\cdot (\theta_0). \]
\end{prop}

\begin{proof}
We know by Proposition \ref{prop-solution-torsor} that the $w_1^2 a^{n-1}$ solutions for a fixed cover $f\colon C\to E$ uniformely spread among the possible correlators, so that we get some integer multiple of the torsor $\vartheta(f)\cdot(\theta_0)$ indexing the possible correlators. To conclude, we merely need to compute the cardinality of $f(\Tor_\delta(C))$.

The basis $(e_1,e_2)$ diagonalizing the lattice inclusion provides real coordinates on $E$ and $C$ such that $f$ has the following expression:
\[
f \colon \ C\simeq (\RR/\ZZ)^2 \longrightarrow E\simeq (\RR/\ZZ)^2, \qquad
 (u,v) \longmapsto (ku,(a/k)v )
 \]
from which we see that \smash{$|f(\Tor_\delta(C))|=\frac{\delta^2}{\gcd(k,\delta)\gcd(a/k,\delta)}$}. Therefore, each correlator carries
\[
a^{n-1}\left(\frac{w_1}{\delta}\right)^2\gcd(k,\delta)\gcd(a/k,\delta)
\]
 solutions. Summing over covers yields the result.
\end{proof}

Our next goal is to find an expression avoiding the summation over the morphisms $f\colon C\to E$, i.e., the index $a$ sublattices of $\pi_1(E)\simeq\ZZ^2$.



\subsection[Computation of the (theta\_0)-coefficient]{Computation of the $\boldsymbol{(\theta_0)}$-coefficient}
\label{sec-0-coeff}

The next step toward the proof of Theorem~\ref{theo-expression-local-correlated-invariant} is to use Proposition \ref{prop-first-expression-local-correlated-inv} to compute the $(\theta_0)$-coefficient, where $\theta_0$ is a special correlator, belonging to the support of all terms in the sum.

To do so, we need to know the number of sublattices of $\pi_1(E)\simeq\ZZ^2$ of a given type, which is precisely the definition of the Dedekind $\psi$-function: it the unique multiplicative function such that for any prime number $p$, $\psi(p^\alpha)=p^{\alpha-1}(p+1)$. There are $\psi(n)$ index $n$ sublattices of $\ZZ^2$ having type $(1,n)$, which are also known as \textit{primitive} sublattices. Merely dividing by $k$, we get that there are $\psi\bigl(\frac{a}{k^2}\bigr)$ sublattices of type $(k,a/k)$.

\begin{expl}
We have $\psi(2)=3$ as there are $3$ sublattices of $\ZZ^2$ of index $2$. If $(e_1,e_2)$ is a~basis of $\ZZ^2$, these lattices are $\langle 2e_1,e_2\rangle$, $\langle e_1,2e_2\rangle$ and $\langle e_1+e_2,e_1-e_2\rangle$.
\end{expl}

\begin{remarkk}
In particular, as there are $\sigma(a)$ index $a$ sublattices, we get that
\[ \sum_{k^2|a}\psi\left(\frac{a}{k^2}\right)=\sigma(a). \]
This may also be checked using the multiplicativity of $\psi$ and $\sigma$ and the fact that the identity is true on powers of prime numbers.
\end{remarkk}

Using Proposition \ref{prop-first-expression-local-correlated-inv}, we get the following expression.

\begin{lem}
Let \smash{$s_\delta(a)=\sum_{k^2|a}\gcd(k,\delta)\gcd(a/k,\delta)\psi\bigl(\frac{a}{k^2}\bigr)$}. For the level $\delta$ refinement, the~$(\theta_0)$-coefficient is equal to $a^{n-1}\bigl(\frac{w}{\delta}\bigr)^2s_\delta(a)$.
\end{lem}

\begin{proof}
We need to compute the $(0)$-coefficient of \smash{$\sum_{f\colon C\to E}\gcd(k,\delta)\gcd(a/k,\delta)\vartheta(f)$}. As each~$\vartheta$ has by definition coefficient $1$ at $(0)$, and as there are $\psi\bigl(\frac{a}{k^2}\bigr)$ lattices of type $(k,a/k)$, we get the result.
\end{proof}

Before getting to the computation of the full local correlated invariant, we transform the above expression to better suit for our purposes, transforming the sum over $k^2|a$ into a sum over~$d|\delta$ using arithmetic functions more convenient than $\psi$.
\begin{itemize}\itemsep=0pt %[label=$\triangleright$]
 \item We already introduced the sum of divisors function $\sigma(a)=\sum_{d|a}d$, whose Dirichlet generating series is $\zeta(s)\zeta(s-1)$.
 \item Given $d\in\NN$, the Dirichlet generating series of $\overline{\sigma}^d(a)=\mathds{1}_{d|a}\sigma(a/d)$ is
\[ \sum_{a=1}^\infty \frac{\overline{\sigma}^d(a)}{a^s} = \sum_{a'=1}^\infty \frac{\sigma(a')}{(da')^s}=\frac{1}{d^s}\zeta(s)\zeta(s-1). \]
 \item The Dirichlet generating series of the second Jordan function $J_2$ is
\[ \sum_{d=1}^\infty \frac{J_2(d)}{d^s}=\frac{\zeta(s-2)}{\zeta(s)}. \]
\end{itemize}

\begin{lem}
The function $s_\delta(a)$ is fully multiplicative in the following sense:
\[ s_\delta(a)=\prod s_{p^{\nu_p(\delta)}}\bigl(p^{\nu_p(a)}\bigr), \]
where products are over the set of prime numbers. Furthermore, we have the following expression:
\[ s_\delta(a)=\sum_{d|\delta} J_2(d)\overline{\sigma}^d(a). \]
\end{lem}

\begin{proof}
We first prove the multiplicativity: assume that we have the decomposition in products of primes $a=\prod p^{\alpha_p}$ and $\delta=\prod p^{\delta_p}$. The sum over $k^2|a$ writes itself as sum over the families~$(i_p)$ with $2i_p\leqslant \alpha_p$ indexed by the set of prime numbers $\P$
\[ \sum_{k^2|a}\gcd(k,\delta)\gcd(a/k,\delta)\psi\bigl(a/k^2\bigr) = \sum_{2i_p\leqslant \alpha_p}\prod \psi\bigl(p^{\alpha_p-2i_p}\bigr)p^{\min(i_p,\delta_p)+\min(\alpha_p-i_p,\delta_p)}, \]
using the multiplicativity of the Dedekind function $\psi$. We can factor the sum as a product and get the sought multiplicativity.

To prove the desired identity, we compute both Dirichlet generating series. The multiplicativity allows to factor the generating series as a product over prime numbers so that
\[ \sum_{a,\delta=1}^\infty \frac{s_\delta(a)}{\delta^sa^t} = \prod_{p\in\P} \biggl(\sum_{d,2i\leqslant\alpha} \psi\bigl(p^{\alpha-2i}\bigr)p^{\min(i,d)+\min(\alpha-i,d)}p^{-sd-t\alpha} \biggr). \]
To finish the computation of the Dirichlet series for $s_\delta(a)$, the inner sum may be rewritten
\[ \sum_{i,j,k}\psi\bigl(p^i\bigr)p^{\min(j,k)+\min(i+j,k)}p^{-ks-2jt-it}, \]
where $p$ is a prime number and $i,j,k\geqslant 0$. Setting aside $i=0$, $\psi\bigl(p^i\bigr)=p^{i-1}(p+1)$ and we can split according to the value of $k$ and compute the geometric sums, yielding some rational function in $p^{-t}$ and $p^{-s}$.

For the second expression, the Dirichlet series is
\begin{align*}
\sum_{\delta,a} \frac{\sum_{d|\delta} J_2(d)\overline{\sigma}^d(a)}{\delta^s a^t} ={} & \sum_{d,k,a}\frac{J_2(d)\overline{\sigma}^d(a)}{d^sk^sa^t} \qquad(\text{where}\quad\delta=dk) \\
= {}& \sum_{d,k} \frac{J_2(d)}{d^s}\frac{1}{k^s}\frac{1}{d^t}\zeta(t)\zeta(t-1)
= \frac{\zeta(s+t-2)}{\zeta(s+t)}\zeta(s)\zeta(t)\zeta(t-1).
\end{align*}
Using that $\zeta(s)=\prod\frac{1}{1-p^{-s}}$, it can also be expressed as the product of values of a rational function at prime numbers. By a technical but elementary computation, we check that both rational functions are actually the same, so that generating functions have equal coefficients, finishing the proof of the identity.
\end{proof}

Notice that if $\gcd(a,\delta)=1$, we have $\overline{\sigma}^d(a)=0$ for every divisor $d$ of $\delta$. Then, we have $s_\delta(a)=\sigma(a)$. This way, the remaining terms may appear as a correction term when $\gcd(a,\delta)\neq 1$.



\begin{expl}
If $p$, $q$ are distinct prime numbers, we have the following values:
 \begin{gather*}
s_p = \sigma+\bigl(p^2-1\bigr)\overline{\sigma}^p,\qquad
s_{p^2}= \sigma+\bigl(p^2-1\bigr)\overline{\sigma}^p+\bigl(p^4-p^2\bigr)\overline{\sigma}^{p^2},\\
s_{pq} = \sigma + \bigl(p^2-1\bigr)\overline{\sigma}^p + \bigl(q^2-1\bigr)\overline{\sigma}^q+\bigl(p^2-1\bigr)\bigl(q^2-1\bigr)\overline{\sigma}^{pq}.
\end{gather*}
\end{expl}


\subsection{Case of other correlators}

To finish the proof of Theorem~\ref{theo-expression-local-correlated-invariant}, we now study the case where $\theta\neq\theta_0$ is another correlator. We proceed in several steps:
 \begin{itemize}\itemsep=0pt %[label=$\circ$]
 \item First prove that the coefficients do not depend on the precise correlator $\theta$ but only on the order of $\theta-\theta_0$ inside $\Tor_\delta(E)$. To do so, we use the deformation invariance for a suitable family where we deform the base curve $E$.
 \item Then, through Lemma \ref{lem-unrefinement}, we prove an induction relation that enables a formal computation of the correlated invariant for $\theta\neq \theta_0$.
 \item Using the induction relation, we show the multiplicativity of the coefficient functions, so that we may restrict to the case of powers of primes, which we compute also using the induction. Combination of both leads to Theorem~\ref{theo-expression-local-correlated-invariant}.
 \end{itemize}

\subsubsection{Dependence on the order} Assumed that $L=\O$ and choose $\theta_0=0$, so that the torsor of correlators is actually $\Tor_\delta(E)$. For $\theta\in\Tor_\delta(E)$, its order $\omega(\theta)$ is the smallest positive integer $n$ such that $n\theta\equiv 0$.

\begin{lem}
 Assuming $L=\O$, the correlated invariant \smash{$\langle \pt_0,1_{w_1},\pt_{w_2},\dots,\pt_{w_n}\rangle_{1,a[E],\bfw}^\theta$} only depends on the order $\omega(\theta)$.
\end{lem}

\begin{proof}
 Write $E$ as some $E_\tau=\CC/\langle1;\tau\rangle$ for some complex number $\tau$ in the Poincar\'e half-plane. For any \smash{$\gamma=\bigl(\begin{smallmatrix}
 a & b \\ c & d \\
 \end{smallmatrix}\bigr)\in {\rm SL}_2(\ZZ)$}, we have an isomorphism between $E_\tau$ and $E_{\gamma\cdot\tau}$ where \smash{$\gamma\cdot\tau=\frac{a\tau+b}{c\tau+d}$}. The isomorphism is actually given by
 \[ \varphi\colon\ (z \modulo 1,\tau) \longmapsto \frac{z}{c\tau+d} \quad \modulo 1,\gamma\cdot\tau. \]
 Choose a path $\tau(t)$ in the Poincar\'e half-plane going from $\tau(0)=\tau$ to $\tau(1)=\gamma\cdot\tau$. The $\delta$-torsion elements in $E_{\tau(t)}$ are deformed continuously along with $\tau(t)$ as they can be written as $u\tau(t)+v$, where $u,v\in\Tor_\delta(\RR/\ZZ)$. The identification $\varphi\colon E_\tau\to E_{\gamma\cdot\tau}$ then induces some monodromy among torsion elements. More precisely, if $u,v\in \Tor_\delta(\RR/\ZZ)$, we have
 \[ \varphi^{-1}\left(u\frac{a\tau+b}{c\tau+d}+v \right)=u(a\tau+b)+v(c\tau+d) = (au+cv)\tau+(bu+dv). \]
 Therefore, we deduce that the values of the correlated invariants for the torsion elements $\theta=u\tau+v$ and $\gamma\cdot\theta=(au+cv)\tau+(bu+dv)$ are the same. As the action of ${\rm SL}_2(\ZZ)$ on the set of torsion elements of the same order in $\Tor_\delta\bigl((\RR/\ZZ)^2\bigr)$ is transitive, we conclude.
\end{proof}

\subsubsection{Coefficients, induction relation and multiplicativity} For $r|\delta$, let $s_\delta[r](a)$ be such that for any $\theta\in\Tor_\delta(E)$,
\[ \langle\pt_0,1_w,\pt_{w_2},\dots,\pt_{w_n}\rangle_{1,a[E],\bfw}^\theta = a^{n-1}\left(\frac{w}{\delta}\right)^2 s_\delta[\omega(\theta)](a). \]
In particular, we have $s_\delta[1]=s_\delta$, corresponding to $\theta=0$. We also set the global function with values in the group algebra \smash{$S_\delta(a) = \sum_{\delta\theta\equiv 0} s_\delta[\omega(\theta)](a)\cdot (\theta)$}, so that
\[ \langle\langle\pt_0,1_w,\pt_{w_2},\dots,\pt_{w_n}\rangle\rangle^\delta_{1,a[E],\bfw} = a^{n-1}\left(\frac{w}{\delta}\right)^2 S_\delta(a). \]
The definition of $S_\delta$ directly comes from the invariants. Theorem~\ref{theo-expression-local-correlated-invariant} consists in proving that $\frac{S_\delta}{\delta^2}=\boldsymbol{\sigma}_\delta$, as defined before Theorem~\ref{theo-expression-local-correlated-invariant}.

\begin{lem}\label{lem-relation-sdelta}
For every pair $\delta'|\delta$, we have the following relation:
\[ \sum_{r|\delta'}J_2(r)s_\delta[r](a) = (\delta')^2s_{\delta/\delta'}(a). \]
\end{lem}

\begin{proof}
We already have $\m{\delta'}\bigl(\fvir{\M(a,\bfw)}{\delta}\bigr)=\fvir{\M(a,\bfw)}{\delta/\delta'}$. Integrating the point constraints and looking at the $(0)$-coefficient yields
\[ a^{n-1}\left(\frac{w}{\delta/\delta'}\right)^2 s_{\delta/\delta'}(a) = \sum_{\delta'\theta\equiv 0}a^{n-1}\left(\frac{w}{\delta}\right)^2s_\delta[\omega(\theta)](a). \]
The sum is over $\theta$ such that $\delta'\theta\equiv 0$. For any $r|\delta'$, there are $J_2(r)$ elements of order exactly $r$, yielding the desired relation.
\end{proof}

These relations form a system which is triangular for the order given by the divisibility. They are thus sufficient to compute all the functions $s_\delta[r]$. We carry out some examples below.

\begin{expl}
If $\delta=p$ is a prime number, we have
$s_p[1]+\bigl(p^2-1\bigr)s_p[p] =p^2\sigma$.
As we already know that $s_p[1]=s_p=\sigma+\bigl(p^2-1\bigr)\overline{\sigma}^p$, we deduce that
\[ \bigl(p^2-1\bigr)\overline{\sigma}^p+\bigl(p^2-1\bigr)s_p[p] = \bigl(p^2-1\bigr)\sigma, \]
and thus
$s_p[p]=\sigma-\overline{\sigma}^p$.
\end{expl}

\begin{expl}
For $\delta=p^2$, we have the following equations:
\begin{gather*}
 s_{p^2}+\bigl(p^2-1\bigr)s_{p^2}[p]+\bigl(p^4-p^2\bigr)s_{p^2}\big[p^2\big] = p^4s_1 =p^4\sigma, \\
 s_{p^2}+\bigl(p^2-1\bigr)s_{p^2}[p] = p^2s_p =p^2\sigma+\bigl(p^4-p^2\bigr)\overline{\sigma}^p, \\
 s_{p^2} = s_{p^2} = \sigma+\bigl(p^2-1\bigr)\overline{\sigma}^p+\bigl(p^4-p^2\bigr)\overline{\sigma}^{p^2},
 \end{gather*}
 which solves for
\begin{gather*}
 s_{p^2} = \sigma+\bigl(p^2-1\bigr)\overline{\sigma}^p+\bigl(p^4-p^2\bigr)\overline{\sigma}^{p^2}, \qquad
 s_{p^2}[p] = \sigma+\bigl(p^2-1\bigr)\overline{\sigma}^p -p^2\overline{\sigma}^{p^2}, \\
 s_{p^2}[p^2] = \sigma-\overline{\sigma}^p.
 \end{gather*}
\end{expl}

To reduce the computation down to the powers of primes, we use the induction relation to prove the multiplicativity of the functions $s_\delta[r](a)$ and $S_\delta(a)$.

\begin{prop}\label{prop-multiplicativity-S}
The function $s_\delta[r](a)$ and $S_\delta(a)$ are fully multiplicative: for pairs of coprime elements $(r_1,r_2)$, $(a_1,a_2)$ and $(\delta_1,\delta_2)$, we have
\[ s_{\delta_1\delta_2}[r_1r_2](a_1a_2) = s_{\delta_1}[r_1](a_1)s_{\delta_2}[r_2](a_2)\qquad
\text{and}\qquad
S_{\delta_1\delta_2}(a_1a_2)=S_{\delta_1}(a_1)S_{\delta_2}(a_2). \]
\end{prop}

\begin{proof}
We show multiplicativity by induction. It is true for $r=1$ by multiplicativity of $s_\delta$. For the induction step, we use the identity from Lemma \ref{lem-relation-sdelta} and the multiplicativity of $s_\delta$,
\begin{gather*}
(\delta'_1)^2(\delta'_2)^2 s_{\delta_1/\delta'_1}(a_1)s_{\delta_2/\delta'_2}(a_2) = (\delta'_1\delta'_2)^2 s_{\delta_1\delta_2/\delta'_1\delta'_2}(a_1a_2), \\
\sum_{\substack{r_1|\delta'_1 \\ r_2|\delta'_2}} J_2(r_1)J_2(r_2)s_{\delta_1}[r_1](a_1)s_{\delta_2}[r_2](a_2) = \sum_{r|\delta'_1\delta'_2} J_2(r)s_{\delta_1\delta_2}[r](a_1a_2) \\
 \phantom{\sum_{\substack{r_1|\delta'_1 \\ r_2|\delta'_2}} J_2(r_1)J_2(r_2)s_{\delta_1}[r_1](a_1)s_{\delta_2}[r_2](a_2) }{}=\sum_{\substack{r_1|\delta'_1 \\ r_2|\delta'_2}} J_2(r_1)J_2(r_2)s_{\delta_1\delta_2}[r_1r_2](a_1a_2).
\end{gather*}
To get the last sum, we use that each divisor of $\delta'_1\delta'_2$ can uniquely been written as a product of a divisor of $\delta'_1$ and a divisor of $\delta'_2$. If we assume multiplicativity for $r$ a strict divisor of $\delta'_1\delta'_2$, all the terms in the sum except the last one are equal, and we are left with the multiplicativity for~${r=\delta'_1\delta'_2}$, finishing the induction. The multiplicativity of $S_\delta$ follows from the multiplicativity of~$s_\delta[r]$,
\begin{gather*}
\biggl(\sum_{\delta_1\theta_1\equiv 0}s_{\delta_1}[\omega(\theta_1)](a_1)\cdot(\theta_1) \biggr)\biggl(\sum_{\delta_2\theta_2\equiv 0}s_{\delta_2}[\omega(\theta_2)](a_2)\cdot(\theta_2) \biggr) \\
\qquad= \sum_{\delta_1\theta_1\equiv\delta_2\theta_2\equiv 0}s_{\delta_1}[\omega(\theta_1)](a_1)s_{\delta_2}[\omega(\theta_2)](a_2)\cdot(\theta_1+\theta_2).
\end{gather*}
As each $\delta_1\delta_2$ torsion element can be written uniquely as the sum of a $\delta_1$-torsion and a $\delta_2$-torsion element, we conclude.
\end{proof}

\subsubsection{Computation for powers of primes} Because of multiplicativity, we only need to compute the function when $a,r,\delta$ are the power of a common prime $p$. Recall that we defined the elements of $\QQ[E]$: $\vartheta_d=\frac{1}{d^2}\sum_{d\theta\equiv 0}(\theta)$. We also momentarily set \smash{$\vartheta_d^{\mathrm{prim}}=\sum_{\omega(\theta)=d}(\theta)$}.


\begin{prop}\label{prop-expression-powers-primes}
 The function $(\delta,r,a)\mapsto s_\delta[r](a)$ has the following values over powers of primes, for $r>0$
 \[ s_{p^d}[p^r]=s_{p^{d-r}}-p^{2d-2r}\overline{\sigma}^{p^{d-r+1}} = \sum_0^{d-r} J_2\bigl(p^j\bigr)\overline{\sigma}^{p^j} -p^{2d-2r}\overline{\sigma}^{p^{d-r+1}} . \]
 Consequently, we find the following expressions for $S_{p^d}$:
 \[ \frac{S_{p^d}}{p^{2d}} = \sum_{r=0}^{d-1} \bigl(\overline{\sigma}^{p^r}-\overline{\sigma}^{p^{r+1}}\bigr)\vartheta_{p^{d-r}} + \overline{\sigma}^{p^d}\vartheta_1
 = \sigma\vartheta_{p^d} + \sum_1^d \overline{\sigma}^{p^r} (\vartheta_{p^{d-r}}-\vartheta_{p^{d-r+1}}). \]

\end{prop}

\begin{proof}
 By Lemma \ref{lem-relation-sdelta}, we have the following relations:
 \[ p^{2r}s_{p^{d-r}} = \sum_{j=0}^{r} J_2\bigl(p^j\bigr) s_{p^d}\big[p^j\big]. \]
 For $r>0$, making the difference between relations for $r$ and $r-1$, we get
 \[ p^{2r}s_{p^{d-r}}-p^{2r-2}s_{p^{d-r+1}} = J_2(p^r)s_{p^d}[p^r] = p^{2r-2}\bigl(p^2-1\bigr)s_{p^d}[p^r]. \]
 Therefore, after dividing by $p^{2r-2}$, using that \smash{$s_{p^{d-r+1}}=s_{p^{d-r}}+J_2\bigl(p^{d-r+1}\bigr)\overline{\sigma}^{p^{d-r+1}}$}, we finally get
 \[ \bigl(p^2-1\bigr)s_{p^{d-r}} -J_2\bigl(p^{d-r+1}\bigr)\overline{\sigma}^{p^{d-r+1}} = \bigl(p^2-1\bigr)s_{p^d}[p^r]. \]
 Dividing by $p^2-1$ yields the desired expression. Then, we deduce that $S_{p^d}$ has the following expression:
 \begin{align*}
 S_{p^d} ={} & s_{p^d}\vartheta_1 + \sum_{r=1}^d \bigl(s_{p^{d-r}}-p^{2(d-r)}\overline{\sigma}^{p^{d-r+1}} \bigr)\cdot \vartheta_{p^r}^\mathrm{prim} \\
 ={} & s_{p^d}\vartheta_1 + \sum_{r=1}^d \bigl(s_{p^{d-r}}-p^{2(d-r)}\overline{\sigma}^{p^{d-r+1}} \bigr)\cdot \bigl(p^{2r}\vartheta_{p^r}-p^{2r-2}\vartheta_{p^{r-1}} \bigr).
 \end{align*}
 We can now split the sum performing a summation by part to get
 \begin{align*}
 S_{p^d}= {}& \bigl(s_{p^d}-s_{p^{d-1}}+p^{2(d-1)}\overline{\sigma}^{p^d}\bigr)\vartheta_1 +\bigl(\sigma-\overline{\sigma}^p\bigr)p^{2d}\vartheta_{p^d} \\
 & + \sum_1^{d-1} p^{2r}\vartheta_{p^r} \bigl(s_{p^{d-r}} -p^{2(d-r)}\overline{\sigma}^{p^{d-r+1}}-s_{p^{d-r-1}}+p^{2(d-r-1)}\overline{\sigma}^{p^{d-r}} \bigr) \\
 ={} & p^{2d}\overline{\sigma}^{p^d}\vartheta_1 +p^{2d}\sum_1^d \bigl(\overline{\sigma}^{p^{d-r}} - \overline{\sigma}^{p^{d-r+1}} \bigr)\vartheta_{p^r}.
 \end{align*}
 This yields the first expression. Performing a second summation by part yields the second expression.
\end{proof}

We have two expressions of $S_{p^d}$, depending on which feature we wish to emphasize: the $\vartheta_d$ or the $\overline{\sigma}^d$. The first ones are useful to compute in the group algebra since they are projectors, while the second are multiplicative functions that provide quasi-modularity properties.

\begin{proof}[Proof of Theorem~\ref{theo-expression-local-correlated-invariant}]
Combining the multiplicativity from Proposition \ref{prop-multiplicativity-S} and the expression for powers of primes from Proposition \ref{prop-expression-powers-primes} yields and explicit expression for $S_\delta$ and we finally get that $\frac{S_\delta}{\delta^2}=\boldsymbol{\sigma}_\delta$.
\end{proof}
